\let\negmedspace\undefined
\let\negthickspace\undefined
\documentclass[journal]{IEEEtran}
\usepackage[a5paper, margin=10mm, onecolumn]{geometry}
\usepackage{tfrupee}

\setlength{\headheight}{1cm}
\setlength{\headsep}{0mm}

\usepackage{gvv-book}
\usepackage{cite}
\usepackage{amsmath,amssymb,amsfonts,amsthm}
\usepackage{algorithmic}
\usepackage{graphicx}
\usepackage{textcomp}
\usepackage{xcolor}
\usepackage{txfonts}
\usepackage{listings}
\usepackage{enumitem}
\usepackage{mathtools}
\usepackage{gensymb}
\usepackage{comment}
\usepackage[breaklinks=true]{hyperref}
\usepackage{tkz-euclide}
\usepackage{listings}
\usepackage{gvv}
\usepackage{tikz}
\def\inputGnumericTable{}
\usepackage[latin1]{inputenc}
\usepackage{color}
\usepackage{array}
\usepackage{longtable}
\usepackage{calc}
\usepackage{multirow}
\usepackage{hhline}
\usepackage{ifthen}
\usepackage{lscape}
\begin{document}

\bibliographystyle{IEEEtran}

\title{4.10.18}
\author{Taraka Abhinav - EE25BTECH11016}
{\let\newpage\relax\maketitle}

\renewcommand{\thefigure}{\theenumi}
\renewcommand{\thetable}{\theenumi}
\setlength{\intextsep}{10pt}

\numberwithin{figure}{enumi}
\renewcommand{\thetable}{\theenumi}

\textbf{Question}:\\
Find the area of the region bounded by the lines $y=4x+5$, $y=5-x$, and $4y=x+5$.

\solution \\
The bounded region is a triangle. Its area can be determined by first finding the coordinates of its three vertices, which are the points of intersection of the given lines.\\

\textbf{Step 1: Calculation of Vertices}\\

\textbf{Vertex A:} Intersection of $L_1: -4x+y=5$ and $L_2: x+y=5$
\begin{align}
\left( \begin{array}{cc|c} -4 & 1 & 5 \\ 1 & 1 & 5 \end{array} \right) \xrightarrow{4\rho_2+\rho_1} \left( \begin{array}{cc|c} -4 & 1 & 5 \\ 0 & 5 & 25 \end{array} \right)
\end{align}
From the second row, $5y=25 \implies y=5$. Substituting into the first row, $-4x+5=5 \implies x=0$.
The coordinates for Vertex A are $\myvec{0 \\ 5}$.

\textbf{Vertex B:} Intersection of $L_2: x+y=5$ and $L_3: -x+4y=5$
\begin{align}
\left( \begin{array}{cc|c} 1 & 1 & 5 \\ -1 & 4 & 5 \end{array} \right) \xrightarrow{\rho_1+\rho_2} \left( \begin{array}{cc|c} 1 & 1 & 5 \\ 0 & 5 & 10 \end{array} \right)
\end{align}
From the second row, $5y=10 \implies y=2$. Substituting into the first row, $x+2=5 \implies x=3$.
The coordinates for Vertex B are $\myvec{3 \\ 2}$.

\textbf{Vertex C:} Intersection of $L_1: -4x+y=5$ and $L_3: -x+4y=5$
\begin{align}
\left( \begin{array}{cc|c} -4 & 1 & 5 \\ -1 & 4 & 5 \end{array} \right) \xrightarrow{-4\rho_2+\rho_1} \left( \begin{array}{cc|c} -4 & 1 & 5 \\ 0 & -15 & -15 \end{array} \right)
\end{align}
From the second row, $-15y=-15 \implies y=1$. Substituting into the first row, $-4x+1=5 \implies x=-1$.
The coordinates for Vertex C are $\myvec{-1 \\ 1}$.

\textbf{Step 2: Calculation of Area}\\
The area of a triangle with vertices $\vec{A}$, $\vec{B}$, and $\vec{C}$ can be calculated using the determinant of the vectors forming two of its sides. The formula is:
\[
\text{Area} = \frac{1}{2} \left| \det(\vec{A}-\vec{B}, \vec{A}-\vec{C}) \right|
\]
The given vertices are represented by 2D vectors:
\[
\vec{A} = \myvec{0 \\ 5}, \quad \vec{B} = \myvec{3 \\ 2}, \quad \vec{C} = \myvec{-1 \\ 1}
\]
First, we find the two side vectors:
\begin{align}
    \vec{A}-\vec{B} &= \myvec{0 \\ 5} - \myvec{3 \\ 2} = \myvec{-3 \\ 3} \label{eq:vec_ab} \\
    \vec{A}-\vec{C} &= \myvec{0 \\ 5} - \myvec{-1 \\ 1} = \myvec{1 \\ 4} \label{eq:vec_ac}
\end{align}
Next, we calculate the determinant of the matrix formed by these two vectors as columns:
\begin{align}
    \det(\vec{A}-\vec{B}, \vec{A}-\vec{C}) &=  \mydet{
\begin{pmatrix}
-3 & 1 \\
3 & 4
\end{pmatrix}
} \nonumber \\
    &= (-3)(4) - (1)(3) \nonumber \\
    &= -12 - 3 = -15 \label{eq:determinant}
\end{align}
The area is half the absolute value of this determinant:
\begin{align}
    \text{Area} = \frac{1}{2} |-15| = 7.5
\end{align}
Thus, the area of the triangle is $7.5$ square units.
\begin{figure}[H]
    \centering
    \includegraphics[width=1\linewidth]{figs/Figure_1.png}
    \label{fig:placeholder}
\end{figure}
\end{document}